\section{Related Works}
\label{sec:related_works}

\subsection{Hallucination in Vision-Language Models: Detection and Benchmarks}

Hallucination in large Vision-Language Models (VLMs) has emerged as one of the most pressing reliability concerns in multimodal AI. Early systematic work by \citet{li2023pope} formalized the problem of object hallucination through the POPE benchmark, which probes model responses to binary existence questions under random, popular, and adversarial sampling conditions. This evaluation protocol remains a standard reference for assessing object-level faithfulness and is one of the primary benchmarks used to evaluate CMVKG-Guard. Complementing POPE at a higher level of complexity, \citet{sun2024aligning} introduced MMHal-Bench, a diagnostic suite of 96 image--question pairs spanning eight question categories and twelve object topics, designed to expose both perceptual and factual hallucinations in instruction-following VLMs. Their analysis of failure modes under factually augmented RLHF provides a taxonomic grounding for the three-layer verification hierarchy --- visual alignment, knowledge grounding, and logical reasoning --- adopted in our Cross-Modal Verification Engine. Broader surveys on hallucination in multimodal large language models \citep{bai2024hallucination} further characterize hallucination along perceptual, factual, and reasoning axes, motivating the multi-granular approach central to CMVKG-Guard.

\subsection{Training-Free Hallucination Mitigation}

A significant line of research seeks to reduce hallucinations at inference time without modifying model weights, which directly motivates the training-free design of CMVKG-Guard. Visual Contrastive Decoding (VCD) \citep{leng2024mitigating} contrasts output distributions produced from original and visually perturbed inputs, suppressing language-prior-driven hallucinations without any auxiliary training. OPERA \citep{huang2024opera} identifies that hallucinations correlate with excessive attention concentration on summary tokens and addresses this through an over-trust penalty and retrospective allocation during beam search. Both methods establish strong inference-time baselines and are direct comparators for CMVKG-Guard in our experimental evaluation.

More recent 2025 work has introduced even more sophisticated training-free mechanisms. DeGF \citep{zhang2025self} employs text-to-image generative feedback as a cross-modal verification signal: a candidate VLM response is used to synthesize a reference image, which is then diffed against the original input to identify perceptual discrepancies and trigger corrective re-decoding. MARINE \citep{zhao2025mitigating}, awarded ICML 2025 Spotlight recognition, extracts structured object-level information from open-source vision expert models and injects it as image-grounded guidance during token generation, outperforming even fine-tuned baselines on POPE and CHAIR. Whereas DeGF and MARINE each employ a single cross-modal feedback channel, CMVKG-Guard consolidates visual alignment, knowledge grounding, and logical reasoning into a unified three-layer verification engine operating over a dynamically constructed knowledge graph.

\subsection{Post-Hoc Hallucination Correction and Self-Verification}

An orthogonal strategy corrects hallucinated outputs after generation rather than modifying the decoding process. Woodpecker \citep{yin2024woodpecker} implements a five-stage pipeline that extracts key concepts from a generated response, formulates targeted visual questions, validates claims using specialist vision models, and rewrites the response based on grounded evidence. LURE \citep{zhou2024analyzing} takes a statistical approach, training a lightweight post-hoc revisor conditioned on three identified hallucination factors --- co-occurrence bias, generation uncertainty, and spatial position --- achieving a 23\% improvement on CHAIR metrics. The most recent self-verification framework, REVERSE \citep{wu2025generate}, unifies generation adjustment with on-the-fly verification by introducing special confidence tokens and retrospective resampling, achieving state-of-the-art hallucination reduction of up to 34\% on HaloQuest. Unlike these post-hoc methods, CMVKG-Guard operates as an end-to-end middleware that intercepts and corrects token candidates in real time during inference, with a measured overhead of less than 15\,ms per token, making it suitable for latency-sensitive deployment.

\subsection{Multimodal Knowledge Graph Construction and Reasoning}

Knowledge graphs have long served as structured repositories for encoding relational world knowledge, and recent work has extended this paradigm to multimodal settings. MR-MKG \citep{lee2024multimodal} demonstrates that multimodal knowledge graphs, incorporating both visual and textual entity relations, substantially improve performance on multimodal question answering and analogy reasoning when used as auxiliary context for LLMs, with strong gains achievable by training only 2.25\% of model parameters. VaLiK \citep{liu2025aligning}, presented at ICCV 2025, advances this direction by introducing the first fully annotation-free and zero-shot pipeline for constructing MMKGs from raw images through cascaded VLMs, augmented with cross-modal similarity filtering to suppress noise. The Dynamic Hierarchical Multimodal Knowledge Graph (DH-MMKG) builder in CMVKG-Guard extends these ideas by constructing scene graphs on-the-fly during inference and enriching them with external knowledge without requiring any offline MMKG pre-construction, enabling real-time adaptability to arbitrary input queries.

\subsection{Knowledge Graph-Enhanced Retrieval-Augmented Generation for VLMs}

Retrieval-Augmented Generation (RAG) has proven effective at grounding language model outputs in verifiable factual knowledge, and recent work has adapted this paradigm to the multimodal setting. mKG-RAG \citep{wang2024mkgrag} constructs multimodal knowledge graphs through MLLM-powered keyword extraction and vision-text matching, then applies a dual-stage retrieval strategy to supply structured context for visual question answering, achieving a 20\% absolute improvement over zero-shot baselines on InfoSeek and E-VQA. MMGraphRAG \citep{wan2025mmgraphrag} further integrates scene-graph-based visual content refinement, spectral clustering for cross-modal entity linking, and knowledge-path-guided retrieval into a unified RAG framework, reporting state-of-the-art results on DocBench and MMLongBench. CMVKG-Guard shares the core insight of both frameworks --- that structured multimodal relational knowledge substantially reduces hallucination in knowledge-intensive tasks --- but goes beyond retrieval by incorporating active verification and adaptive correction at the token level. The knowledge-grounded factuality improvements achieved by the RULE framework \citep{xia2024rule} in medical VLMs further validate the viability of RAG-based factuality control in high-stakes domains, consistent with CMVKG-Guard's design goals for healthcare deployment.

\subsection{Scene Graph Generation}

Scene graphs provide structured, relational representations of visual content and serve as a critical intermediate representation for grounding VLM outputs in perceptual evidence. Panoptic video scene graph generation \citep{yang2023panoptic} extended the scene graph paradigm to the temporal domain, demonstrating that pixel-level segmentation masks can ground dynamic relational structures in video sequences with high fidelity. This work informs the on-the-fly scene graph construction component of DH-MMKG, which builds hierarchical entity-relation graphs from individual frames during inference. The utility of scene graph representations for hallucination mitigation is directly demonstrated by \citet{wu2025combating}, who show that grounding VLM reasoning in bottom-up scene graph structures substantially reduces object, attribute, and relationship hallucinations by anchoring the reasoning chain in perceptually verified relational facts, consistent with the visual alignment layer in CMVKG-Guard's verification engine.

\subsection{Confidence Calibration in Large Language and Vision-Language Models}

Reliable uncertainty quantification is a prerequisite for adaptive, risk-sensitive inference. \citet{xiong2024llms} provides the first systematic empirical evaluation of confidence elicitation strategies for black-box LLMs, benchmarking prompting strategies, sampling methods, and aggregation techniques across five models on calibration and failure-prediction metrics. Their findings establish that well-designed verbalized confidence scores can be meaningfully calibrated and are predictive of factual correctness. Building on this, \citet{zhang2024calibrating} decomposes LLM confidence into two orthogonal components—uncertainty about the question and fidelity to the predicted answer—and propose UF Calibration. This plug-and-play method achieves robust calibration across temperature settings. The Adaptive Confidence Calibration mechanism in CMVKG-Guard operationalizes these insights by dynamically adjusting detection thresholds based on both query complexity and domain risk, applying higher verification stringency in safety-critical contexts while maintaining low latency overhead in routine inference.